% =============================================================================
%   III. DATA AND METHODOLOGY
% =============================================================================
\section{Data and Methodology}
\label{sec:data}

\subsection{Data Sources and Sample Construction}
\label{sec:datacollect}

The analysis uses a cleaned cross-sectional sample of \textbf{276 mid-cap equities}
as of March~2025, spanning \textbf{186 U.S.} and \textbf{90 Indian} firms
(from S\&P MidCap 400, Russell Midcap, and NIFTY Midcap 150 ESG constituent
sets) \cite{nse_esg}.
The final sample spans 11 GICS sectors: Healthcare,
Consumer Cyclical, Industrials, Technology, Financial Services, Basic
Materials, Consumer Defensive, Real Estate, Communication Services, Energy,
and Utilities.

Primary inputs are aggregated from four source families. First, Yahoo Finance
provides financial statement fields, market-history series, and governance risk
scores (e.g., audit, board, compensation, shareholder-rights, and overall risk)
\cite{yahoo_finance, yfinance}, with governance coverage for 213/276 firms.
Second, SEC EDGAR XBRL filings provide observed governance- and disclosure-linked
fields for 102/276 firms \cite{sec_edgar}. Third, financial-proxy ESG
construction is used only where direct observations are unavailable (6/276 firms
at the firm-level coverage frontier). Fourth, benchmark index series are used for
validation windows: 753 trading days for the S\&P 500 and
739 trading days for the NIFTY 50 over the 2023--2026 period.

All Indian monetary fields are converted to USD using a fixed
INR/USD rate of 83.0 (March~2025 reference rate), ensuring cross-country
comparability and reproducibility of all ratio constructions.

\subsubsection{ESG Data: Hybrid Proxy Pipeline}
\label{sec:proxy_esg}

Because open-source ESG coverage is sparse and uneven for mid-caps, ESG inputs
are constructed through a hybrid provenance hierarchy, applied at each
\textit{company $\times$ indicator} cell:
\begin{enumerate}
  \item \textbf{Real Yahoo Finance governance risk scores} (213/276 firms),
  \item \textbf{Real SEC EDGAR governance disclosures} (102/276 firms),
  \item \textbf{Financial-proxy-derived indicators} (6/276 firms requiring fallback),
  \item \textbf{Sector-median imputation},
  \item \textbf{Cross-sector (global) imputation},
  \item \textbf{Residual missing value (NaN)}.
\end{enumerate}

The ESG module contains 12 indicators across Environmental (E), Social (S),
and Governance (G) pillars. A central transparency constraint is that \textbf{E
and S are predominantly proxy-derived from financial/accounting signals},
whereas \textbf{G has substantial directly observed data} via Yahoo governance
risk and SEC governance-related disclosures. Examples include R\&D intensity as
an innovation/workforce-development proxy and energy/cost-structure-derived
signals as environmental-intensity proxies \cite{konar2001does, edmans_2011,
adams2009women, bebchuk2009what}.

Three binary ESG indicators (carbon reduction target, human rights policy,
anti-corruption policy) exhibit zero variance post-imputation (all
companies~$=1.0$). Pre-imputation coverage for these indicators was
approximately 15--20\%, with sector-median imputation defaulting to 1.0 in
sectors where the majority of companies with available data reported having
such policies. These indicators are retained for completeness but contribute
zero discriminating power to the ESG pillar scores; the scoring engine
automatically excludes constant-valued indicators from the weighted average
(see Section~\ref{sec:scoring}).

Every ESG value is written with explicit provenance metadata
(\texttt{real\_sec}, \texttt{real\_yahoo}, \texttt{financial\_proxy},
\texttt{sector\_imputed}, \texttt{global\_imputed}, \texttt{missing}),
enabling exact reproducibility and post-hoc sensitivity tests by source tier.

\begin{figure}[t]
\centering
\includegraphics[width=\columnwidth]{fig17_missing_data}
\caption{Missing-data structure before and after preprocessing in the
276-company panel. The ESG block is the dominant source of sparsity pre-cleaning;
coverage improves through tiered proxying and imputation with provenance tracking.}
\label{fig:missing_data}
\end{figure}

\subsection{Variable Construction and Treatment}
\label{sec:vartypes}

After schema harmonization and cleaning, the analytical matrix contains
\textbf{202 variables}. Outlier treatment uses adaptive winsorization at the
1st/99th percentiles, clipping 530 observations in total.

For distributional stabilization, 15 highly skewed columns are log-transformed
prior to normalization. Core standardization uses a robust median/MAD transform:
\begin{equation}
z_{ij}^{\mathrm{MAD}} = \frac{x_{ij} - \mathrm{median}(x_j)}{\mathrm{MAD}(x_j)}
\end{equation}
followed by score mapping to a bounded index scale:
\begin{equation}
\mathrm{score}_{ij} = \min\{100,\max[0,\,50 + 20\,z_{ij}^{\mathrm{MAD}}]\}.
\end{equation}

Table~\ref{tab:vartypes} summarizes the variable-type framework used for
cleaning and transformations.

\begin{table}[t]
\centering
\caption{Variable-Type Classification and Treatment}
\label{tab:vartypes}
\small
\setlength{\tabcolsep}{3pt}
\begin{tabularx}{\columnwidth}{@{}l>{\raggedright\arraybackslash}p{2.1cm}>{\raggedright\arraybackslash}X@{}}
\toprule
\textbf{Type} & \textbf{Examples} & \textbf{Treatment} \\
\midrule
Ratio & leverage, margins, turnover & 1/99 winsorization; MAD-$z$ normalization \\
Ordinal & disclosure level scales & rank-preserving transform; median fill \\
Binary & ESG commitments (SBTi, UNGC, CDP, RE100) & mode fill; consistency checks (0/1) \\
Categorical & sector, country tags & one-hot or groupwise conditioning only \\
\bottomrule
\end{tabularx}
\end{table}

\begin{figure}[t]
\centering
\includegraphics[width=\columnwidth]{fig01_score_distributions}
\caption{Cross-sectional distributions of the nine factor scores after robust
normalization and clipping to $[0,100]$.}
\label{fig:score_distributions}
\end{figure}

\subsection{Outlier Detection and Robust Scaling}
\label{sec:outlier}

Outlier detection follows a consensus rule across IQR fences,
modified MAD-$z$, and classical $z$-scores \cite{tukey_1977, iglewicz_1993,
aguinis_2013}. Values are winsorized (not deleted) to preserve sector balance
and avoid reducing already sparse ESG coverage.

\subsection{Nine-Factor Scoring Model and Investor Profiles}
\label{sec:scoring}

The composite model contains 9 factors:
\begin{enumerate}
  \item ESG Composite (baseline pillar weights E=20\%, S=35\%, G=45\%, with sector-specific SASB E/S/G weighting),
  \item Financial Score (profitability 30\%, scale 20\%, efficiency 20\%, stability 15\%, valuation 15\%),
  \item Market Score (liquidity 40\%, volatility 30\%, momentum 30\%)\cite{amihud_2002},
  \item Operational Score (productivity, innovation),
  \item Risk-Adjusted Score (Sharpe, Sortino, drawdown),
  \item Value Score (inverted P/E, P/B, P/S, and EV multiples),
  \item Growth Score (revenue and earnings growth),
  \item Stability Score (liquidity 35\%, leverage 35\%, volatility 30\%),
  \item Sector Position (within-sector percentile).
\end{enumerate}

The \texttt{similarity\_rank} factor was removed from the deployed composite
after collinearity diagnostics showed high overlap with
\texttt{ESG\_composite} ($r=0.70$). The deployed model therefore uses 9 factors after excluding \texttt{similarity\_rank} due to this collinearity. Its prior profile budget was redistributed
across the remaining nine factors while preserving the profile-level total
weight budget. Variance-inflation diagnostics remain low after deployment
($\max\,\mathrm{VIF}=2.80$).

Governance inputs from Yahoo and SEC are harmonized before G-pillar
construction using source-wise quantile normalization. This removes provider
scale effects before merge and mitigates the observed cross-provider
divergence ($r=0.59$).

For valuation signals, lower raw multiples imply higher attractiveness, so all
value ratios are direction-inverted before standardization. In periods where
the value factor exhibits persistently negative IC (as in the 2020--2024
growth-led regime), the deployed value signal is flipped to capture the
observed growth premium while preserving rank monotonicity.

The stability block is configured in three categories: liquidity (35\%), using
\texttt{current\_ratio} and \texttt{quick\_ratio} at 50/50 weights; leverage
(35\%), using inverse \texttt{debt\_to\_equity} (35\%), inverse
\texttt{debt\_to\_ebitda} (30\%), and \texttt{cash\_flow\_to\_debt} (35\%);
and volatility (30\%), using inverse \texttt{price\_volatility\_30d}.

For firm $i$ and profile $p$, the composite score is
\begin{equation}
C_i^{(p)} = \sum_{k=1}^{9} w_k^{(p)}\,F_{ik}
\label{eq:composite}
\end{equation}
with $\sum_k w_k^{(p)}=1$. To prevent circularity in forward-return validation,
the default scoring profile excludes \texttt{market\_score} from deployment-time
weights and re-normalizes the remaining eight-factor vector.

\subsubsection{Re-standardisation and Rank Aggregation}
\label{sec:rank_aggregation}

Before profile aggregation, each factor is re-standardized cross-sectionally:
\begin{equation}
\tilde{F}_{ik} = 50 + 10\cdot\frac{F_{ik}-\bar{F}_{\cdot k}}{\sigma(F_{\cdot k})}
\label{eq:restandardise}
\end{equation}
and optionally rank-mapped for percentile-based aggregation:
\begin{equation}
R_{ik} = 100\cdot\frac{\mathrm{rank}(\tilde{F}_{ik})}{N}.
\label{eq:rank_transform}
\end{equation}

\begin{algorithm}[t]
\caption{Composite Scoring with Provenance-Aware ESG Inputs}
\label{alg:composite}
\begin{algorithmic}[1]
\REQUIRE Cleaned feature matrix $X$, ESG provenance tensor $P$, sector map $s(i)$, profile weights $w^{(p)}$
\ENSURE Composite scores $C_i^{(p)}$ and ranked lists by profile
\STATE Build ESG indicators via six-tier hierarchy; retain source tag in $P$
\STATE Apply 1/99 winsorization and log-transform selected skewed variables
\STATE Compute robust normalized indicators using MAD-$z$
\STATE Construct 9 factor scores $F_{ik}$ from indicator blocks
\STATE Compute sector-percentile position
\STATE Re-standardize factors via Eq.~\ref{eq:restandardise}
\FOR{each investor profile $p$}
  \STATE If circularity-controlled mode: set $w_{\text{market}}^{(p)}\leftarrow0$ and renormalize
  \STATE Compute $C_i^{(p)}$ using Eq.~\ref{eq:composite}
\ENDFOR
\STATE Return sorted rankings, factor attributions, and provenance audit table
\end{algorithmic}
\end{algorithm}

\begin{table}[t]
\centering
\caption{Factor Weights (\%) for Three Investor Profiles}
\label{tab:profile_weights}
\small
\begin{tabular}{@{}lrrr@{}}
\toprule
\textbf{Factor} & \textbf{ESG-First} & \textbf{Balanced} & \textbf{Financial-First} \\
\midrule
ESG Composite      & 55 & 22 &  5 \\
Financial Score     &  8 & 18 & 25 \\
Market Score        &  2 &  5 & 10 \\
Operational Score   &  5 &  6 &  8 \\
Risk-Adjusted Score &  8 & 15 & 18 \\
Growth Score        &  6 & 14 & 15 \\
Value Score         &  3 &  4 &  7 \\
Stability Score     &  6 &  8 &  7 \\
Sector Position     &  7 &  8 &  5 \\
\midrule
\textbf{Total}      & 100 & 100 & 100 \\
\bottomrule
\end{tabular}
\end{table}

\subsection{SASB Materiality Weighting}
\label{sec:sasb_materiality}

Sector-specific ESG pillar weights are based on the SASB Materiality Map,
which links industry context to financially material ESG channels
\cite{khan_2016, sasb_2023}. For each sector $s$,
$w_E^{(s)} + w_S^{(s)} + w_G^{(s)} = 1$ and
$\mathrm{ESG}_i^{(s)} = w_E^{(s)}E_i + w_S^{(s)}S_i + w_G^{(s)}G_i$.

\begin{table}[t]
\centering
\caption{SASB-Informed E/S/G Pillar Weights by Sector}
\label{tab:sasb_weights}
\small
\begin{tabular}{@{}lrrr@{}}
\toprule
\textbf{Sector} & \textbf{E} & \textbf{S} & \textbf{G} \\
\midrule
Energy                 & 0.55 & 0.20 & 0.25 \\
Utilities              & 0.50 & 0.20 & 0.30 \\
Basic Materials        & 0.50 & 0.25 & 0.25 \\
Industrials            & 0.40 & 0.30 & 0.30 \\
Consumer Cyclical      & 0.25 & 0.40 & 0.35 \\
Consumer Defensive     & 0.30 & 0.35 & 0.35 \\
Healthcare             & 0.25 & 0.40 & 0.35 \\
Technology             & 0.20 & 0.45 & 0.35 \\
Financial Services     & 0.15 & 0.35 & 0.50 \\
Real Estate            & 0.40 & 0.25 & 0.35 \\
Communication Services & 0.20 & 0.45 & 0.35 \\
\bottomrule
\end{tabular}
\end{table}

\subsection{Circularity Correction and Validation Design}
\label{sec:circularity}

Circularity is controlled in three ways: (i) excluding market-derived factors
from default deployment weights, (ii) reporting validation metrics for both
market-inclusive and ex-market composites, and (iii) decoupling
$E\_\text{score}$ from \texttt{financial\_score} via regression residualization.
Specifically, we estimate
\begin{equation}
E_i = \alpha + \beta\,\mathrm{financial\_score}_i + \varepsilon_i,
\end{equation}
and use residualized environmental signal $E_i^{\perp}=\hat{\varepsilon}_i$ in
sensitivity checks to reduce construction circularity.

\subsection{Statistical Validation Framework}
\label{sec:weight_method}
\label{sec:factor_validity}

Validation combines seven complementary procedures:
\begin{enumerate}
  \item \textbf{Bootstrap stability:} 500 iterations of profile-score recomputation \cite{efron_1993}.
  \item \textbf{Subsampling stability:} 500 replications using 80\% random subsamples.
  \item \textbf{Stratified $k$-fold cross-validation:} folds are stratified by sector to preserve cross-sector composition.
  \item \textbf{Leave-one-sector-out cross-validation:} each sector is held out in turn to test transfer to unseen industries; mean test IC across held-out sectors is 0.200.
  \item \textbf{Constrained weight estimation:} L2 regularization ($\lambda=0.25$) with an entropy penalty, under per-factor bounds $w_k\in[0.02,0.35]$.
  \item \textbf{Walk-forward validation:} 165 out-of-sample train/test splits.
  \item \textbf{Factor ablation and regime analysis:} remove each factor $k$ and evaluate bull/bear partitions.
\end{enumerate}

To further eliminate forward-looking contamination from momentum, the return
validation target is a forward quality proxy,
\texttt{forward\_quality\_proxy} (formerly
\texttt{composite\_forward\_return}), built from forward ROA, ROE, net margin,
revenue growth, operating margin, and current ratio. This decouples validation
from momentum, which remains inside \texttt{market\_score} construction.

Monotonicity correction is applied before final rank aggregation. Each factor is
first blended using 70\% cross-sectional and 30\% within-sector scoring,
followed by continuous isotonic regression on 10 decile bins. The isotonic
blend weight is adaptive (60--90\%), with higher isotonic weight assigned to
factors exhibiting weaker pre-correction monotonicity.

Capacity analysis is performed under mid-cap execution constraints: a 10-day
execution window, 3\% average daily volume participation, and a 15\% maximum
single-name position cap as a fraction of ADV. Trading impact is modeled following the simplified square-root specification of Almgren and Chriss \cite{almgren_chriss_2001},
\begin{equation}
\mathrm{impact\_bps} = 10\times\sqrt{\frac{\mathrm{position\_value}}{\mathrm{ADV}}},
\end{equation}
and feasible deployment is stress-tested at \$25M, \$50M, \$100M, \$250M,
\$500M, \$1B, and \$2B AUM levels.

Turnover robustness is evaluated with score perturbations of $\pm10\%$,
$\pm20\%$, and $\pm30\%$, each with 200 Monte Carlo simulations.

Generalization is evaluated by transferring the framework to a 45-stock
large-cap benchmark panel. Cross-profile consistency under transfer is assessed
with Friedman rank tests and Kendall's $W$ concordance.

Rank stability across settings is summarized with Kendall rank correlation:
\begin{equation}
\tau = \frac{n_c - n_d}{\tfrac{1}{2}N(N-1)}
\end{equation}
where $n_c$ and $n_d$ are concordant and discordant pair counts.

Cross-sectional predictive content is measured using the Information
Coefficient (IC):
\begin{equation}
\mathrm{IC}_t = \mathrm{corr}\!\left(\mathrm{rank}(C_{i,t}),\;r_{i,t+h}\right)
\end{equation}
with horizon $h\in\{1,3,6\}$ months.

Cross-sectional asset-pricing validation additionally uses Fama--MacBeth
regressions with sector fixed effects and a size control,
$\log(\mathrm{market\_cap})$, to isolate incremental factor premia.

The cross-sectional information ratio aggregates IC through time:
\begin{equation}
\mathrm{IR}_{\mathrm{CS}} = \frac{\mathbb{E}[\mathrm{IC}_t]}{\mathrm{Std}(\mathrm{IC}_t)}.
\end{equation}

ESG is also evaluated as a risk filter. Firms are sorted into ESG-score
quintiles and compared on realized volatility, maximum drawdown, beta,
skewness, and Sharpe ratio. Complementing quintile contrasts, Spearman
correlations are computed between ESG scores and each risk metric to quantify
direction and strength of risk alignment.

\begin{figure}[t]
\centering
\includegraphics[width=\columnwidth]{fig10_weight_sensitivity}
\caption{Weight-perturbation stability from 500 bootstrap trials and 80\%
subsampling. Error bars denote 95\% confidence intervals of rank outcomes.}
\label{fig:weight_sensitivity}
\end{figure}

\paragraph{Power Analysis}
\label{sec:power_analysis}
With $N=276$ firms in cross-section and repeated validation folds,
the design has adequate power for moderate effect sizes in rank-correlation
tests while preserving sector-level heterogeneity. Specifically, for a
Spearman correlation test at $\alpha=0.05$ (two-sided) with $N=276$, the
minimum detectable effect size is $\rho \approx 0.11$ at power $= 0.80$,
which is sufficient for detecting economically meaningful factor-return
associations. The bootstrap and subsampling procedures ($B=500$)
further ensure adequate resolution for confidence-interval estimation
\cite{efron_1993}.
